% GENERATED FILE. Edit canonical lessons, then run npm run generate:books.
% canonical-source-hash: fnv1a64:765f94e62be6f5ed
% canonical-lessons: GU-C538-vimasan, GU-C538-dhacupcu, GU-C538-gadamthal, GU-C538-gholvu, GU-C538-nivedo

\chapter{In Two Minds}
\label{ch:gu-in-two-minds}
\hlchaptermodality{538}

\begin{chapteropening}
\textbf{By the end of this chapter:} \emph{I can say that I can't make up my mind and why.}

Complete the last lesson of chapter 538: I can say that I can't make up my mind and why.
\end{chapteropening}

\section[\texorpdfstring{\gu{વિમાસણ} (vimāsaṇ)}{વિમાસણ (vimāsaṇ)}]{\gu{વિમાસણ} (vimāsaṇ) — a quandary}

\label{lesson:GU-C538-vimasan}



\begin{quote}
Before the new one: say the Gujarati for to pick out, then the Gujarati for to appeal.
\end{quote}

\subsection*{You'll want to know: \gu{વિમાસણ}}
\textbf{\gu{વિમાસણ}} — \emph{vimāsaṇ} — \textquotedblleft{}a quandary\textquotedblright{}.

\begin{quote}
\gu{જો} \gu{તમે} \gu{વિમાસણમાં} \gu{પડશો} \gu{તો} \gu{મને} \gu{ફોન} \gu{કરો}.
\end{quote}

\begin{quote}
\emph{jo tame vimāsaṇmā̃ paḍśo to mane phon karo.}
\end{quote}

\begin{quote}
If you find yourself in a quandary, phone me.
\end{quote}

A feminine noun.

\begin{cousinweb}
\textbf{\gu{વિમાસણ}} (\emph{vimāsaṇ}) is being caught between choices and uneasy about what will come; one \textquotedblleft{}falls\textquotedblright{} into it: \textbf{\gu{વિમાસણમાં} \gu{પડવું}}. CDIAL lists the Gujarati word, with the sense \textquotedblleft{}fearful anticipation\textquotedblright{}, under Sanskrit \emph{vimarśana}, \textquotedblleft{}investigation\textquotedblright{} (CDIAL 11820).
\end{cousinweb}

\begin{grammarlens}[title={What you've built}]
One more word for saying what happens if.
\end{grammarlens}

\subsection*{Your turn}
\begin{itemize}
\raggedright
  \item \emph{Say it:} \emph{vimāsaṇ}
  \item \emph{Say it:} \emph{vimāsaṇ}, once more
  \item \emph{Say it:} \emph{jo tame vimāsaṇmā̃ paḍśo to mane phon karo.} — If you find yourself in a quandary, phone me.
  \item \emph{Recall:} say the Gujarati for to pick out, then the Gujarati for to appeal, then say \emph{vimāsaṇ} again
\end{itemize}

\begin{tcolorbox}[breakable,colback=teal!4,colframe=teal!35!black,title={Before you move on}]
What does \gu{વિમાસણ} mean? (\textquotedblleft{}A quandary\textquotedblright{}.) Say it once more.
\end{tcolorbox}
\section[\texorpdfstring{\gu{ઢચુપચુ} (ḍhacupcu)}{ઢચુપચુ (ḍhacupcu)}]{\gu{ઢચુપચુ} (ḍhacupcu) — undecided (in two minds)}

\label{lesson:GU-C538-dhacupcu}



\begin{quote}
Before the new one: say the Gujarati for to appeal, then the Gujarati for a quandary.
\end{quote}

\subsection*{You'll want to know: \gu{ઢચુપચુ}}
\textbf{\gu{ઢચુપચુ}} — \emph{ḍhacupcu} — \textquotedblleft{}undecided (in two minds)\textquotedblright{}.

\begin{quote}
\gu{જો} \gu{તમે} \gu{ઢચુપચુ} \gu{રહેશો} \gu{તો} \gu{કોઈ} \gu{તમારો} \gu{વિશ્વાસ} \gu{નહીં} \gu{કરે}.
\end{quote}

\begin{quote}
\emph{jo tame ḍhacupcu raheśo to koī tamāro viśvās nahī̃ kare.}
\end{quote}

\begin{quote}
If you stay undecided, nobody will trust you.
\end{quote}

\begin{cousinweb}
\textbf{\gu{ઢચુપચુ}} (\emph{ḍhacupcu}) is an expressive, jingling word, its sound swinging like the mind it describes. No dictionary we use records an older form, so its history is not settled.
\end{cousinweb}

\begin{grammarlens}[title={What you've built}]
One more word for saying what happens if.
\end{grammarlens}

\subsection*{Your turn}
\begin{itemize}
\raggedright
  \item \emph{Say it:} \emph{ḍhacupcu}
  \item \emph{Say it:} \emph{ḍhacupcu}, once more
  \item \emph{Say it:} \emph{jo tame ḍhacupcu raheśo to koī tamāro viśvās nahī̃ kare.} — If you stay undecided, nobody will trust you.
  \item \emph{Recall:} say the Gujarati for to appeal, then the Gujarati for a quandary, then say \emph{ḍhacupcu} again
\end{itemize}

\begin{tcolorbox}[breakable,colback=teal!4,colframe=teal!35!black,title={Before you move on}]
What does \gu{ઢચુપચુ} mean? (\textquotedblleft{}Undecided (in two minds)\textquotedblright{}.) Say it once more.
\end{tcolorbox}
\section[\texorpdfstring{\gu{ગડમથલ} (gaḍamthal)}{ગડમથલ (gaḍamthal)}]{\gu{ગડમથલ} (gaḍamthal) — inner turmoil}

\label{lesson:GU-C538-gadamthal}



\begin{quote}
Before the new one: say the Gujarati for a quandary, then the Gujarati for undecided.
\end{quote}

\subsection*{You'll want to know: \gu{ગડમથલ}}
\textbf{\gu{ગડમથલ}} — \emph{gaḍamthal} — \textquotedblleft{}inner turmoil\textquotedblright{}.

\begin{quote}
\gu{જો} \gu{હું} \gu{તેને} \gu{સાચું} \gu{કહેત} \gu{તો} \gu{મારી} \gu{ગડમથલ} \gu{ઓછી} \gu{થાત}.
\end{quote}

\begin{quote}
\emph{jo hũ tene sācũ kahet to mārī gaḍamthal ochī thāt.}
\end{quote}

\begin{quote}
If I had told him the truth, my inner turmoil would have been less.
\end{quote}

A feminine noun.

\begin{cousinweb}
\textbf{\gu{ગડમથલ}} (\emph{gaḍamthal}) is an expressive word for the churning of a mind that cannot settle. No dictionary we use records an older form, so its history is not settled.
\end{cousinweb}

\begin{grammarlens}[title={What you've built}]
One more word for saying what happens if.
\end{grammarlens}

\subsection*{Your turn}
\begin{itemize}
\raggedright
  \item \emph{Say it:} \emph{gaḍamthal}
  \item \emph{Say it:} \emph{gaḍamthal}, once more
  \item \emph{Say it:} \emph{jo hũ tene sācũ kahet to mārī gaḍamthal ochī thāt.} — If I had told him the truth, my inner turmoil would have been less.
  \item \emph{Recall:} say the Gujarati for a quandary, then the Gujarati for undecided, then say \emph{gaḍamthal} again
\end{itemize}

\begin{tcolorbox}[breakable,colback=teal!4,colframe=teal!35!black,title={Before you move on}]
What does \gu{ગડમથલ} mean? (\textquotedblleft{}Inner turmoil\textquotedblright{}.) Say it once more.
\end{tcolorbox}
\section[\texorpdfstring{\gu{ઘોળવું} (ghoḷvũ)}{ઘોળવું (ghoḷvũ)}]{\gu{ઘોળવું} (ghoḷvũ) — to mull over}

\label{lesson:GU-C538-gholvu}



\begin{quote}
Before the new one: say the Gujarati for undecided, then the Gujarati for inner turmoil.
\end{quote}

\subsection*{You'll want to know: \gu{ઘોળવું}}
\textbf{\gu{ઘોળવું}} — \emph{ghoḷvũ} — \textquotedblleft{}to mull over\textquotedblright{}.

\begin{quote}
\gu{જો} \gu{તમે} \gu{આ} \gu{વાત} \gu{બે} \gu{દિવસ} \gu{ઘોળશો} \gu{તો} \gu{તમને} \gu{જવાબ} \gu{મળશે}.
\end{quote}

\begin{quote}
\emph{jo tame ā vāt be divas ghoḷśo to tamne javāb maḷśe.}
\end{quote}

\begin{quote}
If you mull this over for two days, you'll find the answer.
\end{quote}

Past, with \gu{તેણે} (he or she …): \gu{તેણે} \gu{વાત} \gu{ઘોળી}, \emph{teṇe vāt ghoḷī}.

\begin{cousinweb}
\textbf{\gu{ઘોળવું}} (\emph{ghoḷvũ}) is filed by CDIAL under Sanskrit \emph{ghōlayati}, \textquotedblleft{}stirs, mixes\textquotedblright{}, and CDIAL lists the Gujarati verb, \textquotedblleft{}to shake about; to ponder on\textquotedblright{} (CDIAL 4526): a thought stirred round and round.
\end{cousinweb}

\begin{grammarlens}[title={What you've built}]
One more word for saying what happens if.
\end{grammarlens}

\subsection*{Your turn}
\begin{itemize}
\raggedright
  \item \emph{Say it:} \emph{ghoḷvũ}
  \item \emph{Say it:} \emph{ghoḷvũ}, once more
  \item \emph{Say it:} \emph{jo tame ā vāt be divas ghoḷśo to tamne javāb maḷśe.} — If you mull this over for two days, you'll find the answer.
  \item \emph{Recall:} say the Gujarati for undecided, then the Gujarati for inner turmoil, then say \emph{ghoḷvũ} again
\end{itemize}

\begin{tcolorbox}[breakable,colback=teal!4,colframe=teal!35!black,title={Before you move on}]
What does \gu{ઘોળવું} mean? (\textquotedblleft{}To mull over\textquotedblright{}.) Say it once more.
\end{tcolorbox}
\section[\texorpdfstring{\gu{નિવેડો} (niveḍo)}{નિવેડો (niveḍo)}]{\gu{નિવેડો} (niveḍo) — a settlement (of a matter)}

\label{lesson:GU-C538-nivedo}



\begin{quote}
Before the new one: say the Gujarati for inner turmoil, then the Gujarati for to mull over.
\end{quote}

\subsection*{You'll want to know: \gu{નિવેડો}}
\textbf{\gu{નિવેડો}} — \emph{niveḍo} — \textquotedblleft{}a settlement (of a matter)\textquotedblright{}.

\begin{quote}
\gu{જો} \gu{આપણે} \gu{બેસીને} \gu{વાત} \gu{કરીશું} \gu{તો} \gu{નિવેડો} \gu{આવશે}.
\end{quote}

\begin{quote}
\emph{jo āpaṇe besīne vāt karīśũ to niveḍo āvśe.}
\end{quote}

\begin{quote}
If we sit down and talk, the matter will be settled.
\end{quote}

A masculine noun.

\begin{cousinweb}
\textbf{\gu{નિવેડો}} (\emph{niveḍo}) is how a matter ends once it is decided; a settlement \textquotedblleft{}comes\textquotedblright{}: \textbf{\gu{નિવેડો} \gu{આવશે}}. CDIAL records the Gujarati noun, \textquotedblleft{}a settlement, a decision\textquotedblright{}, under a reconstructed \emph{*nirvarati}, \textquotedblleft{}restrains\textquotedblright{} (CDIAL 7392).
\end{cousinweb}

\begin{grammarlens}[title={What you've built}]
One more word for saying what happens if.
\end{grammarlens}

\subsection*{Your turn}
\begin{itemize}
\raggedright
  \item \emph{Say it:} \emph{niveḍo}
  \item \emph{Say it:} \emph{niveḍo}, once more
  \item \emph{Say it:} \emph{jo āpaṇe besīne vāt karīśũ to niveḍo āvśe.} — If we sit down and talk, the matter will be settled.
  \item \emph{Recall:} say the Gujarati for inner turmoil, then the Gujarati for to mull over, then say \emph{niveḍo} again
\end{itemize}

\begin{tcolorbox}[breakable,colback=teal!4,colframe=teal!35!black,title={Before you move on}]
What does \gu{નિવેડો} mean? (\textquotedblleft{}A settlement (of a matter)\textquotedblright{}.) Say it once more.
\end{tcolorbox}
